% --- Sezione Controllo Verticale ---
\section{\hbox{Analisi di Lyapunov per il controllo verticale}}
Di seguito è analizzata la stabilità secondo Lyapunov del sistema d'errore attraverso l'applicazione della legge di controllo \textit{Sliding Mode} sviluppata per il volo verticale. Si considera, a titolo esemplificativo, il caso lungo l'asse $z$.\\
Si adotta un sistema di riferimento \textbf{NED (North-East-Down)}, in cui l'asse $z$ è orientato verso il centro della Terra; pertanto, valori di $z$ positivi indicano una quota crescente verso il basso.

\subsection{Modellazione della Dinamica}
La dinamica del sistema viene analizzata nel frame inerziale, poiché il loop di controllo esterno si basa sull'errore di posizione, che perde di significato nel body frame.\\
Il secondo principio della dinamica, proiettato sull'asse $z$ del sistema inerziale, definisce l'evoluzione del velivolo come:
\begin{equation}
    m\ddot{z} = mg + \cos(\phi)\cos(\theta) F_{\mathrm{aero},z} - \cos(\phi)\cos(\theta) F_{\mathrm{thrust},z} + d(t)
    \label{eq:dynamics_z_final}
\end{equation}
dove:
\begin{itemize}
    \item $m$ è la massa totale del velivolo;
    \item $g$ è l'accelerazione di gravità (positiva nel riferimento NED);
    \item $F_{\mathrm{thrust},z}$ è la componente verticale della spinta, ovvero l'ingresso di controllo;
    \item $F_{\mathrm{aero},z} = -\frac{1}{2}\rho S C_d \operatorname{sgn}(\dot{z})\dot{z}^2$ rappresenta la resistenza aerodinamica lungo l'asse verticale;
    \item $d(t)$ rappresenta le incertezze di modello e i disturbi esterni non modellati.
\end{itemize}
Normalizzando rispetto alla massa, si ottiene la forma affine al controllo:
\begin{equation}
    \ddot{z} = g + \cos(\phi)\cos(\theta) \left(\frac{F_{\mathrm{aero},z}}{m} - \frac{u}{m} \right) + \Delta(t)
\end{equation}
In questa formulazione, $\Delta(t) = d(t)/m$ è definita come \textbf{incertezza matched}, in quanto agisce direttamente sullo stesso canale dell'ingresso di controllo $u = F_{\mathrm{thrust},z}$. Si assume che tale disturbo sia limitato superiormente da una costante nota: \hbox{$|\Delta(t)| \le \Delta_{\max}$}, $\forall t \ge 0$.

\subsection{Definizione della \textit{sliding surface}}
Sia $z_{\mathrm{des}}(t)$ la quota di riferimento desiderata e $e(t) = z_{\mathrm{des}}(t) - z(t)$ l'errore. Per un sistema del secondo ordine, si definisce una \textit{sliding surface} $\sigma(t)$ lineare del primo ordine:
\begin{equation}
    \sigma(t) = \dot{e}(t) + \lambda e(t)
\end{equation}
con $\lambda > 0$ costante di progetto. La condizione di regime $\sigma = 0$ definisce una dinamica vincolata (sliding mode) in cui l'errore decade esponenzialmente a zero. 

\subsection{Sintesi della Legge di Controllo}
Per ricavare la legge di controllo, si deriva rispetto al tempo la \textit{sliding surface} $\sigma(t)$, ottenendo $\dot{\sigma} = \ddot{z}_{\mathrm{des}} - \ddot{z} + \lambda \dot{e}$. Sostituendo l'espressione della dinamica del sistema proiettata lungo $z$, si ricava la dinamica dell'errore sulla superficie:
\begin{equation}
    \dot{\sigma} = \ddot{z}_{\mathrm{des}} - g - \cos(\phi)\cos(\theta)\frac{F_{\mathrm{aero},z}}{m} + \cos(\phi)\cos(\theta)\frac{u}{m} - \Delta(t) + \lambda \dot{e}
\end{equation}

L'obiettivo del controllo Sliding Mode è garantire che le traiettorie del sistema raggiungano la superficie $\sigma = 0$ in tempo finito. 
A tal fine, si adotta l'approccio basato sulla \textbf{reaching law}; si impone che la dinamica:
\begin{equation}
    \dot{\sigma} = -k \operatorname{sgn}(\sigma)
\end{equation}
con $k > 0$ parametro di progetto. Uguagliando la dinamica nominale alla legge di raggiungimento desiderata, si ottiene l'equazione algebrica da cui estrarre analiticamente l'ingresso di controllo $u$:
\begin{equation}
    \ddot{z}_{\mathrm{des}} - g - \cos(\phi)\cos(\theta)\frac{F_{\mathrm{aero},z}}{m} + \lambda \dot{e} + \cos(\phi)\cos(\theta)\frac{u}{m} = -k \operatorname{sgn}(\sigma)
\end{equation}

Isolando il termine dipendente dall'azione di controllo, si ottiene:
\begin{equation}
    \cos(\phi)\cos(\theta)\frac{u}{m} = -\ddot{z}_{\mathrm{des}} + g + \cos(\phi)\cos(\theta)\frac{F_{\mathrm{aero},z}}{m} - \lambda \dot{e} - k \operatorname{sgn}(\sigma)
\end{equation}

Si ricava la legge di controllo completa:
\begin{equation}
    u = \frac{m}{\cos(\phi)\cos(\theta)} \left( g - \ddot{z}_{\mathrm{des}} - \lambda \dot{e} \right) + F_{\mathrm{aero},z} - \frac{m}{\cos(\phi)\cos(\theta)} k \operatorname{sgn}(\sigma)
\end{equation}

Affinché l'inversione della dinamica sia matematicamente ben posta, è strettamente necessario introdurre l'assunzione che gli angoli di assetto soddisfino $\phi, \theta \in (-\pi/2, \pi/2)$. Fisicamente, questa condizione evita la singolarità cinematica, garantendo che il velivolo non si trovi in orientamento perfettamente verticale o capovolto, configurazione in cui la spinta non avrebbe alcuna componente utile lungo l'asse $z$ inerziale.

Una parte del controllo annulla la dinamica nota del sistema e assumiamo che i termini residui di queste cancellazioni vengano assorbiti dall'incertezza $\Delta(t)$; un'altra parte del controllo, 
ovvero il contributo discontinuo, impone la legge di raggiungimento.

Infine, applicando la legge di controllo $u$ così sintetizzata al sistema perturbato originale, la dinamica a ciclo chiuso proiettata sulla \textit{sliding surface} si riduce all'espressione:
\begin{equation}
    \dot{\sigma} = -k \operatorname{sgn}(\sigma) - \Delta(t)
\end{equation}
Questo è il punto di partenza per la dimostrazione formale della stabilità, in quanto il parametro $k$ dovrà essere dimensionato opportunamente per dominare il termine incerto $\Delta(t)$.


\subsection{Analisi di Stabilità e Tempo di Raggiungimento}
Per dimostrare la stabilità asintotica e la convergenza in tempo finito alla superficie, si introduce la funzione candidata di Lyapunov $V = \frac{1}{2}\sigma^2$. Dalla sua definizione si evince la relazione:
\begin{equation}
    \sqrt{2V} = |\sigma|
\end{equation}

Calcoliamo la derivata temporale lungo le traiettorie del sistema a ciclo chiuso. Sostituendo la dinamica di $\dot{\sigma}$:
\begin{align}
    \dot{V} &= \sigma \dot{\sigma} = \sigma \left( -k \operatorname{sgn}(\sigma) - \Delta(t) \right) \nonumber \\
            &= -k|\sigma| - \sigma \Delta(t)
\end{align}

Applicando la disuguaglianza triangolare e sapendo che il termine incerto è limitato superiormente ($|\Delta(t)| \le \Delta_{\max}$), si ottiene:
\begin{equation}
    \dot{V} \le -k|\sigma| + \Delta_{\max}|\sigma| = - (k - \Delta_{\max}) |\sigma|
\end{equation}

Definendo $\delta = k - \Delta_{\max}$, e ricordando che $|\sigma| = \sqrt{2V}$, possiamo riscrivere la disequazione differenziale direttamente in termini di $V$:
\begin{equation}
        \dot{V} = \sigma\dot{\sigma} \le -(k - \Delta_{\max})|\sigma| = -\delta \sqrt{2V}
\end{equation}
da cui ricaviamo che il guadagno deve essere $k = \Delta_{\max} + \delta$ per garantire la robustezza del controllo.
\begin{equation}
    \frac{dV}{dt} \le -\delta \sqrt{2V}
\end{equation}

Per dimostrare la convergenza in tempo finito, integriamo la disequazione per separazione di variabili. Separando i termini dipendenti da $V$ da quelli in $t$:
\begin{equation}
    \frac{dV}{\sqrt{V}} \le -\delta \sqrt{2} \, dt
\end{equation}

Integrando ambo i membri tra l'istante iniziale $t=0$ e un generico istante $t$:
\begin{equation}
    \int_{V(0)}^{V(t)} V^{-\frac{1}{2}} \, dV \le -\delta \sqrt{2} \int_{0}^{t} d\tau
\end{equation}

Risolvendo l'integrale si ottiene:
\begin{equation}
    2\sqrt{V(t)} - 2\sqrt{V(0)} \le -\sqrt{2} \, \delta \, t
\end{equation}

Dividendo ambo i membri per $\sqrt{2}$ e reintroducendo la relazione iniziale $\sqrt{2V} = |\sigma|$, l'espressione diventa:
\begin{equation}
    |\sigma(t)| - |\sigma(0)| \le - \delta \, t
\end{equation}

Da cui si definisce l'evoluzione dell'errore nel tempo:
\begin{equation}
    0 \le |\sigma(t)| \le |\sigma(0)| - \delta \, t
\end{equation}

Il tempo di raggiungimento $t_r$ coincide con il momento in cui le traiettorie intersecano per la prima volta la superficie, condizione per la quale si ha $|\sigma(t_r)| = 0$. Sostituendo questa condizione, si ricava il limite superiore per il tempo di convergenza:
\begin{equation}
    t_r \le \frac{|\sigma(0)|}{\delta}
\end{equation}

Questa dimostrazione prova che, scegliendo un guadagno $k$ sufficientemente grande per dominare le incertezze del modello, le traiettorie del sistema vengono forzate a raggiungere la \textit{sliding surface} in un tempo strettamente finito. Una volta sulla superficie ($\sigma=0$), il sistema esibisce la dinamica nominale asintoticamente stabile imposta da $\sigma = \dot{e} + \lambda e = 0$.

\subsection{Mitigazione del Chattering}

La natura discontinua della funzione $\operatorname{sgn}(\sigma)$ induce oscillazioni ad alta frequenza (chattering) che possono eccitare dinamiche non modellate e danneggiare gli attuatori. Per mitigare tale effetto, si introduce la regolarizzazione tramite un \textbf{boundary layer} di spessore $\Phi$, sostituendo la funzione segno con una funzione continua. In questo lavoro si utilizza la tangente iperbolica:
\begin{equation}
u = -k \tanh\left(\frac{\sigma}{\Phi}\right)
\end{equation}

Il parametro $\Phi>0$ definisce lo spessore dello strato limite attorno alla superficie di sliding. In particolare, si introducono le curve di livello
\begin{equation}
\begin{matrix}
L_{\Phi} = \{x : \sigma(x) = \Phi\} \\
L_{-\Phi} = \{x : \sigma(x) = -\Phi\}
\end{matrix}
\end{equation}

che delimitano il \textit{boundary layer}, definito come l'insieme
\begin{equation}
B = \{x \in \mathbb{R}^n : |\sigma(x)| \le \Phi \}.
\end{equation}

L'introduzione dello strato limite comporta la perdita della stabilità asintotica dell'origine. In presenza della regolarizzazione della legge di controllo, il sistema non converge esattamente alla superficie di sliding $\sigma=0$; tuttavia, le traiettorie raggiungono in \textit{tempo finito} un intorno della superficie di sliding e rimangono limitate al suo interno. Per questo motivo si parla di \textbf{stabilità pratica in tempo finito}. In particolare, come riportato in \cite{khalil2002nonlinear}, il sistema risulta \textit{uniformly ultimately bounded}, con un \textit{ultimate bound} la cui ampiezza dipende dal parametro $\Phi$ e dal guadagno del controllore.

Per dimostrare la convergenza verso lo strato limite, si consideri la funzione di Lyapunov candidata
\begin{equation}
V = \frac{1}{2}\sigma^2 .
\end{equation}

Si assuma che la dinamica della variabile di sliding possa essere scritta nella forma
\begin{equation}
\dot{\sigma} = -k \tanh\left(\frac{\sigma}{\Phi}\right) - \Delta(t)
\end{equation}

dove $\Delta(t)$ rappresenta disturbi e incertezze limitate tali che $|\Delta(t)| \le \Delta_{\max}$. La derivata della funzione di Lyapunov lungo le traiettorie del sistema risulta
\begin{equation}
\dot{V} = \sigma \dot{\sigma}
\end{equation}

da cui
\begin{equation}
\dot{V} = -k \sigma \tanh\left(\frac{\sigma}{\Phi}\right) - \sigma \Delta(t).
\end{equation}

Utilizzando il bound sul disturbo si ottiene
\begin{equation}
\dot{V} \le -k |\sigma| \tanh\left(\frac{|\sigma|}{\Phi}\right) + |\sigma| \Delta_{\max}
\end{equation}

ovvero
\begin{equation}
\dot{V} \le -|\sigma|\left(k \tanh\left(\frac{|\sigma|}{\Phi}\right) - \Delta_{\max}\right).
\end{equation}

Al di fuori dello strato limite vale la condizione $|\sigma|>\Phi$, da cui
\[
\frac{|\sigma|}{\Phi} > 1.
\]

Poiché la funzione $\tanh(x)$ è monotona crescente, segue che
\begin{equation}
\tanh\left(\frac{|\sigma|}{\Phi}\right) \ge \tanh(1).
\end{equation}
dove $\tanh(1)$ rappresenta il valore della funzione sul bordo del \textit{boundary layer}.

Pertanto si ottiene la stima
\begin{equation}
\dot{V} \le -|\sigma|\left(k\tanh(1)-\Delta_{\max}\right).
\end{equation}

Definendo $\eta = k \tanh(1) - \Delta_{\max}$, affinché $\dot{V}<0$ al di fuori dello strato limite è sufficiente imporre la seguente condizione sul guadagno del controllore:
\begin{equation}
k = \frac{\Delta_{\max} + \eta}{\tanh(1)}.
\end{equation}

Questa disuguaglianza garantisce che le traiettorie del sistema vengono attratte verso il \textit{boundary layer}.

Poiché $V=\frac{1}{2}\sigma^2$, si ha
\begin{equation}
|\sigma| = \sqrt{2V}.
\end{equation}

Sostituendo tale relazione nella stima precedente si ottiene
\begin{equation}
\dot{V} \le -\eta \sqrt{2V}.
\end{equation}

Questa disuguaglianza differenziale permette di stimare il tempo necessario affinché le traiettorie entrino nello strato limite. Integrando entrambi i membri si ottiene
\begin{equation}
\frac{dV}{\sqrt{V}} \le -\sqrt{2} \eta dt.
\end{equation}

Integrando tra $V(0)$ e $V(t)$ si ricava
\begin{equation}
2(\sqrt{V(t)}-\sqrt{V(0)}) \le -\sqrt{2}\eta t.
\end{equation}

Da cui si ottiene un limite superiore per il tempo necessario affinché la traiettoria raggiunga lo strato limite:
\begin{equation}
t_r \le \frac{|\sigma(0)|}{\eta}.
\end{equation}

Questo risultato dimostra che il sistema converge in \textit{tempo finito} verso il \textit{boundary layer}, confermando la \textbf{finite-time practical stability}.

Una volta entrate nello strato limite, l'argomento della tangente iperbolica è piccolo ($|\sigma|/\Phi \le 1$) e vale l'approssimazione lineare
\begin{equation}
\tanh(x) \approx x.
\end{equation}

La dinamica della variabile di sliding diventa quindi approssimativamente
\begin{equation}
\dot{\sigma} \approx -\frac{k}{\Phi}\sigma - \Delta(t),
\end{equation}

che rappresenta un sistema lineare stabile perturbato. In questa regione le traiettorie rimangono limitate all'interno dello strato limite, con un errore residuo la cui ampiezza dipende dal parametro $\Phi$ e dall'intensità dei disturbi. Riducendo opportunamente il valore di $\Phi$, tale errore può essere reso arbitrariamente piccolo.

In conclusione, emerge un fondamentale trade-off progettuale legato alla scelta dello spessore $\Phi$: un valore ridotto dello strato limite permette di minimizzare l'errore residuo a regime, avvicinando le prestazioni del sistema a quelle dello Sliding Mode ideale; tuttavia, ciò comporta una minore capacità di filtraggio delle oscillazioni ad alta frequenza. Al contrario, un valore di $\Phi$ più elevato garantisce una superiore mitigazione del chattering e una maggiore protezione degli attuatori, accettando però un degradamento della precisione nella stabilità pratica. La taratura del controllore deve quindi bilanciare la fedeltà all'inseguimento della traiettoria con la necessaria robustezza e integrità meccanica del velivolo.

\clearpage